\ncdlesson{14}{Danmark i Europa: Småstatsmentalitet og globalt udsyn}{欧洲中的丹麦：小国心态与全球视野}
{DU3 Modul 5 (PD3) $\rightarrow$ Modul 6 (Studieprøven, B2 $\rightarrow$ C1)｜对标 NCE 4 \& \textit{Historisk-Politisk Essayistik}}
{\textbf{高阶书面连词与关系/副词代词}（原因 \textit{eftersom, idet} · 疑问限定 \textit{hvorvidt, i hvilken udstrækning} · 让步转折 \textit{desuagtet} · 归属列举 \textit{herunder} · 方式结果 \textit{hvormed}）· 历史政论文跨时态衔接与从句语序铁律}
{丹麦在世界舞台上的处世哲学，始终处于一对耐人寻味的张力之中：一方面是 1864 年战败后深植民族记忆的\textbf{“小国心态”（\textit{småstatsmentalitet}）}与对主权让渡的审慎保留；另一方面则是依赖海洋贸易、合作社运动（\textit{andelsbevægelsen}）与绿色外交的\textbf{“全球视野”（\textit{globalt udsyn}）}。本课通过这一历史政论主题，攻克 C1 书面语中最具权威感的七大高阶连词与副词代词。}

% ---------------------------------------------------------------------
\begin{foerduleaeser}
带着下面两个问题阅读 Tekst A 和 Tekst B：
\begin{enumerate}
  \item \textit{Hvordan påvirkede nederlaget i 1864 ifølge Tekst A Danmarks selvforståelse, og hvorfor kan småstatsmentalitet og international åbenhed alligevel eksistere side om side?}（根据 Tekst A，1864 年的战败如何影响了丹麦的自我定位？为什么小国心态与国际开放依然能够并存？）
  \item \textit{Hvilket paradoks peger historikeren i Tekst B på i forbindelse med de danske EU-forbehold og afstemningen i 2022?}（在 Tekst B 中，历史学家围绕丹麦的欧盟保留条款与 2022 年公投指出了哪一悖论？）
\end{enumerate}
\end{foerduleaeser}

% ---------------------------------------------------------------------
\ncdsubsec{二、Tekster（历史政论专篇与学者短评双语篇）}

\subsection*{Tekst A：\textit{Mellem beskyttelse og åbenhed: Danmarks lange vej i Europa}（历史政论专篇 · 约 355 词）}

\begin{tipbox}{语篇说明}
本篇以严谨的历史政论文笔梳理丹麦自 1864 年至今的欧洲定位演变。阅读时请特别关注\textbf{\color{NCDRed}深红色标记的高阶书面连词与关系/副词代词}（\textit{eftersom, idet, hvorvidt, desuagtet, herunder, hvormed, i hvilken udstrækning}）各自引导的语序变化，以及\textbf{\color{NCDNavy}深蓝色标记的政治与历史核心概念}。
\end{tipbox}

For at forstå Danmarks nuværende forhold til Europa må man gå tilbage til det skæbnesvangre \da{nederlag} i 1864, hvor landet mistede omkring en tredjedel af sit territorium til Preussen og Østrig. Med ét slag blev det gamle rige reduceret til en sårbar \da{småstat}, \dabold{hvormed} en dyb forsigtighed over for stormagter slog rod i den politiske kultur. Under slagordet ”Hvad udadtil tabes, skal indadtil vindes” vendte samfundet blikket indad og byggede både folkehøjskoler og en stærk \da{andelsbevægelse}. Denne \da{småstatsmentalitet} \da{prægede} generationer af danskere, \dabold{eftersom} national overlevelse blev kædet tæt sammen med indre sammenhængskraft og folkelig dannelse.

Samtidig har Danmark aldrig haft råd til at isolere sig økonomisk fra omverdenen. Som en lille nation uden store råstoffer har velstanden tværtimod altid hvilet på søfart, landbrugseksport og teknologisk specialisering, \dabold{herunder} medicinalindustri og vedvarende energi. Efter Anden Verdenskrig og især med indtrædelsen i EF i 1973 trådte landet målrettet ind i det europæiske marked, \dabold{idet} en åben og regelbaseret handel skabte trygge rammer for eksporten. Dansk \da{udenrigspolitik} har derfor i årtier balanceret mellem ønsket om fri adgang til globale markeder og frygten for at miste demokratisk kontrol.

Denne indre spænding blev særlig tydelig under folkeafstemningen i 1992, da et snævert flertal af vælgerne sagde nej til Maastrichttraktaten og derved førte til de fire danske \da{EU-forbehold}. Befolkningen diskuterede intenst gennem hele valgkampen, \dabold{hvorvidt} et tættere samarbejde ville udhule lokal \da{selvbestemmelse}, og \dabold{i hvilken udstrækning} en småstat bør \da{afgive} \da{suverænitet} til overnationale organer i Bruxelles. \dabold{Desuagtet} har Danmark i praksis ofte hørt til de medlemslande, der hurtigst og mest loyalt har gennemført fælles europæiske regler på klima-, arbejdsmarkeds- og miljøområdet.

I de seneste år har krigen i Ukraine, ændrede stormagtsrelationer og den globale klimakrise imidlertid tvunget landet til at genoverveje sin traditionelle tilbageholdenhed. Da et klart flertal ved folkeafstemningen i 2022 stemte for at \da{ophæve} forsvarsforbeholdet, markerede det et historisk skift i dansk sikkerhedspolitik. Samtidig søger Danmark i dag aktiv indflydelse gennem grønt diplomati og multilaterale alliancer, \dabold{idet} små stater sjældent kan forme deres egen fremtid alene. Erfaringen viser således, at småstatsmentalitet og \da{globalt udsyn} ikke nødvendigvis er modsætninger, men to uundværlige sider af samme overlevelsesstrategi.

\begin{transbox}
\small
\textbf{\color{NCDTeal}[Tekst A 中文精译]}\\
要理解丹麦如今与欧洲的关系，必须回溯到 1864 年那场改变国运的战败——当时丹麦将大约三分之一的领土割让给了普鲁士与奥地利。转瞬之间，这个古老的王国缩水为一个脆弱的小国，由此一种对大国保持深层警惕的谨慎心态在政治文化中扎下了根。在“凡对外失去的，必于对内赢得”的口号下，丹麦社会将目光转向内部，兴办民众高等学院并建立起强大的合作社运动。这种小国心态深刻塑造了数代丹麦人，因为民族的生存自此与内部凝聚力及民众教养紧密捆绑在了一起。

与此同时，丹麦在经济上从来负担不起与外部世界自我孤立的代价。作为一个缺乏丰富矿产资源的弹丸之地，其繁荣反而始终建立在航海贸易、农产品出口与技术专业化之上——其中包括制药业与可再生能源。二战之后，尤其是 1973 年加入欧洲共同体（EF）以来，丹麦目标明确地融入欧洲市场，因为开放且以规则为基础的贸易为出口创造了安全框架。因此，丹麦外交政策数十年来一直在“渴望自由进入全球市场”与“担忧丧失民主控制权”之间寻求平衡。

这种内在张力在 1992 年公投中表现得尤为明显：当时微弱多数选民对《马斯特里赫特条约》投下反对票，由此催生了丹麦的四项欧盟保留条款。民众在整个竞选期间激烈辩论：更紧密的合作究竟是否会掏空地方自决权，以及一个欧洲小国应当在多大程度上向布鲁塞尔的超国家机构让渡主权。尽管如此，在实践中丹麦却往往跻身于最迅速、最忠实地在气候、劳动力市场与环境领域落实欧盟共同法规的成员国之列。

然而近年来，乌克兰战争、大国关系的变化与全球气候危机迫使这个国家重新审视其传统的克制姿态。当 2022 年公投中明确多数选民投票赞成废除防务保留条款时，这标志着丹麦安全政策的一次历史性转折。与此同时，今天的丹麦正通过绿色外交与多边联盟积极争取影响力，因为小国往往难以凭一己之力塑造自己的未来。历史经验由此表明，小国心态与全球视野未必非此即彼，而是同一套生存战略不可或缺的一体两面。
\end{transbox}

\subsection*{Tekst B：\textit{Kommentar: Når skepsis møder indflydelse i Bruxelles}（历史政治短评 · 约 145 词）}

\textbf{阅读任务（理解与核对）：}\textit{Læs historikeren Thomas Linds kommentar om Danmarks rolle i EU, og besvar spørgsmålene nedenfor.}

\begin{emailbox}{报刊短评｜Fra: Thomas Lind (lektor i europæisk historie) \quad Emne: Småstatens dilemma}
I den danske EU-debat fremstilles \da{suverænitet} ofte som en fast størrelse, man enten beholder udelt derhjemme eller taber i Bruxelles. Denne forsimplede opfattelse overser imidlertid det klassiske dilemma for enhver \da{småstat}: Formelt at stå udenfor garanterer ikke reel \da{selvbestemmelse}, \dabold{eftersom} afgørende beslutninger om sikkerhed, frihandel og grønne produktstandarder alligevel træffes af de store nabolande.

I tre årtier holdt Danmark fast i sine \da{EU-forbehold}, \dabold{herunder} både euro- og forsvarsforbeholdet, ud fra et forståeligt ønske om at beskytte nationale kerneområder. \dabold{Desuagtet} har skiftende danske regeringer måttet følge de fleste europæiske direktiver tæt for ikke at isolere landets eksport og erhvervsliv. Afskaffelsen af forsvarsforbeholdet i 2022 viser tydeligt, \dabold{hvorvidt} udefrakommende kriser kan ændre dybe historiske vaner. Spørgsmålet for fremtiden er derfor ikke længere, om Danmark skal samarbejde med Europa, men \dabold{i hvilken udstrækning} vi selv vil sidde med ved bordet, når de fælles regler skrives.
\end{emailbox}

\textbf{Tekst B 阅读核对（Rigtigt / Forkert / Kort svar，答案见第六部分）：}
\begin{enumerate}
  \item \textit{Ifølge Thomas Lind sikrer det automatisk et lands reelle selvbestemmelse, hvis landet formelt står uden for EU-samarbejdet.} \underline{\hspace{2.2cm}}
  \item \textit{Hvorfor har skiftende danske regeringer alligevel måttet følge de fleste europæiske direktiver trods forbeholdene?} \underline{\hspace{2.2cm}}
  \item \textit{Hvilket afgørende spørgsmål peger Thomas Lind på for Danmarks fremtid?} \underline{\hspace{2.2cm}}
\end{enumerate}

\begin{tipbox}{本课发音与语调要点（Udtale \& Prosodi）}
\begin{itemize}
  \item \textbf{高阶书面连词与代词读音}：\textit{eftersom} \ipa{ˈɛfdɐˌsʌm}（重音在首音节 \textit{ef-}）；\textit{idet} \ipa{iˈdeð} 或弱读 \ipa{iˈdɛ}（注意末尾 \textit{t} 不发硬爆破音，常近于弱化或静音！）；\textit{hvorvidt} \ipa{vɒˈvid}（\textit{h} 与 \textit{d} 均不发音）；\textit{desuagtet} \ipa{dɛsuˈɑɡd̥əð}（重音在 \textit{-ag-}）；\textit{hvormed} \ipa{vɒˈmɛð}（词首 \textit{h} 不发音）。
  \item \textbf{历史政治高频词}：\textit{en suverænitet} \ipa{suvəʁɛniˈtˢeːˀd}（重音在末音节 \textit{-tet}）；\textit{en udstrækning} \ipa{ˈuðˌsdʁɛɡneŋ}（前缀 \textit{ud-} 为软 \textit{d}）；\textit{et forbehold} \ipa{ˈfɒːbeˌhʌlˀ}（后段 \textit{-hold} 带 \textit{stød}，\textit{d} 不发音！）；\textit{et nederlag} \ipa{ˈneːðɐˌlɑʊ̯ˀ}（\textit{-lag} 读滑音 \ipa{lɑʊ̯ˀ}）。
  \item \textbf{长从句政论语调}：遇到 \textit{...eftersom...} 或 \textit{...hvormed...} 引导的后置长从句时，逗号前微升保连贯，切忌在连词后突兀停顿；而在句首 \textit{Desuagtet} 后需紧接动词，一气呵成读出 V2 节奏。
\end{itemize}
\end{tipbox}

\begin{transbox}
\small
\textbf{\color{NCDTeal}[Tekst B 中文对照]}\\
\textbf{学者短评（托马斯·林德，欧洲史副教授｜小国的两难处境）：}在丹麦的欧盟辩论中，主权常常被描绘成一种固定不变的东西——要么原封不动留在国内，要么在布鲁塞尔失去。然而，这种简化的看法忽略了任何一个小国所面临的经典两难：形式上置身事外并不能保证真正的自决权，因为有关安全、自由贸易与绿色产品标准的重大决策终究是由周边大国与邻国做出的。

三十年来，丹麦一直坚持其欧盟保留条款——其中包括欧元与防务保留条款——这是出于一种可以理解的愿望，即保护国家核心领域。尽管如此，历届丹麦政府仍然不得不紧跟大多数欧洲指令，以免孤立本国的出口与工商业界。2022 年废除防务保留条款之举清楚表明，外部危机究竟是否以及如何能够改变根深蒂固的历史习惯。因此，未来的问题不再是丹麦是否应当与欧洲合作，而是在共同规则制定之时，我们在多大程度上愿意亲自坐在谈判桌旁。
\end{transbox}

% ---------------------------------------------------------------------
\ncdsubsec{三、Gloser（核心生词表）}

\begin{tipbox}{词表重点}
本课词汇表囊括了历史政论文与高阶逻辑连接的两大核心集群：政治概念词（如 \textit{suverænitet, selvbestemmelse, forbehold, nederlag}）及七大 C1 书面连接词。特别留意：\textit{et nederlag}（战败/失利）与 \textit{et forbehold}（保留条款）均为单复同形的中性名词；而 \textit{udenrigspolitik} 在加定冠词后缀时须双写 \textit{k}（\textit{udenrigspolitikken}）。
\end{tipbox}

\begin{glosertable}
名词 & en småstat & \textit{(småstaten, småstater, småstaterne)} & 小国，小规模国家 \\
名词 & en småstatsmentalitet & \textit{(småstatsmentaliteten, kun ental)} & 小国心态（谨慎与危机意识） \\
名词 & et udsyn & \textit{(udsynet, udsyn, udsynene)} & 视野，国际格局 \\
名词 & et nederlag & \textit{(nederlaget, nederlag, nederlagene)} & 战败，挫折 \\
名词 & en andelsbevægelse & \textit{(andelsbevægelsen, kun ental)} & 合作社运动（丹麦历史制度） \\
名词 & en udenrigspolitik & \textit{(udenrigspolitikken, -politikker, -erne)} & 外交政策 \\
名词 & et forbehold & \textit{(forbeholdet, forbehold, forbeholdene)} & 保留条款，保留意见 \\
名词 & en suverænitet & \textit{(suveræniteten, kun ental)} & 主权 \\
名词 & en selvbestemmelse & \textit{(selvbestemmelsen, kun ental)} & 自决权，自主决定权 \\
名词 & en udstrækning & \textit{(udstrækningen, udstrækninger, -erne)} & 程度，范围，跨度 \\
\midrule
动词 [har] & at præge & \textit{(præger, prægede, har præget)} & 深刻影响，塑造，打上烙印 \\
动词 [har] & at afgive & \textit{(afgiver, afgav, har afgivet)} & 让渡，交出（权力/主权） \\
动词 [har] & at ophæve & \textit{(ophæver, ophævede, har ophævet)} & 废除，取消（条款/禁令） \\
\midrule
连词 & eftersom & \textit{(从属连词，不变)} & 因为，鉴于（正式因果） \\
连词 & idet & \textit{(从属连词，不变)} & 因为，在于，同时（正式说明） \\
副词 & desuagtet & \textit{(句间副词，不变)} & 尽管如此，即便这样（触发 V2） \\
副词 & herunder & \textit{(副词，不变)} & 其中包括，含以下各项 \\
连词 & hvormed & \textit{(关系副词连词，不变)} & 借此，由此，凭借该举措 \\
\end{glosertable}

% ---------------------------------------------------------------------
\ncdsubsec{四、Grammatik \& Stilistik（高阶书面连词、副词代词与政论时态精讲）}

\begin{grammatikbox}{语法点 1：七大 C1 书面连接词与关系/副词代词全景功能矩阵}
掌握下面这七个高阶书面连接工具，关键在于分清它们的\textbf{句法身份}——究竟是\textbf{从属连词}（引导从句，内部“主语 + 状语 + 动词”不倒装），还是\textbf{句间副词}（位于句首时主句必须发生 V2 倒装）：

\vspace{0.3em}
\renewcommand{\arraystretch}{1.24}
{\small
\begin{tabularx}{\linewidth}{@{} >{\bfseries\raggedright\arraybackslash}p{2.9cm} >{\raggedright\arraybackslash}p{4.6cm} >{\raggedright\arraybackslash}X @{}}
\toprule
\textbf{C1 连接词 / 短语} & \textbf{词性与语序铁律} & \textbf{典型政论示范句} \\
\midrule
\dabold{eftersom}\newline （鉴于，因为） &
\textbf{因果从属连词}：引导客观已知原因；从句内部不倒装，但若放句首，后面主句必倒装。 &
\textit{Mentaliteten prægede landet, \dabold{eftersom} overlevelse krævede sammenhold.} \\
\midrule
\dabold{idet}\newline （由于，在于，同时） &
\textbf{书面说明/同时连词}：常接在主句后补充理据或同步机制；内部绝不倒装。 &
\textit{Landet trådte ind i EF, \dabold{idet} åben handel skabte tryghed.} \\
\midrule
\dabold{hvorvidt}\newline （究竟是否） &
\textbf{间接疑问从属连词}：用于探讨或争论两可议题；内部绝不倒装（禁代以 \textit{hvis}）。 &
\textit{Vælgerne diskuterede, \dabold{hvorvidt} samarbejdet ville udhule demokratiet.} \\
\midrule
\dabold{i hvilken udstrækning}\newline （在多大程度上） &
\textbf{程度疑问短语}：引出定量或界限追问；引导间接问句，内部维持主语 + 动词语序。 &
\textit{Debatten drejer sig om, \dabold{i hvilken udstrækning} man bør afgive suverænitet.} \\
\midrule
\dabold{hvormed}\newline （借此，由此） &
\textbf{关系副词连词}：回指前句整体行动或手段；引导非限定性从句，内部不倒装。 &
\textit{Danmark blev en småstat, \dabold{hvormed} en dyb forsigtighed slog rod.} \\
\midrule
\dabold{desuagtet}\newline （尽管如此） &
\textbf{高阶让步句间副词}（= 更正式的 \textit{alligevel}）：放句首时\textbf{立即触发 V2 倒装}！ &
\textit{\dabold{Desuagtet} har Danmark ofte gennemført fælles regler hurtigt.} \\
\midrule
\dabold{herunder}\newline （其中包括） &
\textbf{归属列举副词}：紧跟名词短语后补充具体子类，不用动词；放句首亦触发 V2。 &
\textit{Eksporten hviler på flere sektorer, \dabold{herunder} vedvarende energi.} \\
\bottomrule
\end{tabularx}
}
\end{grammatikbox}

\begin{grammatikbox}{语法点 2：历史政论文的“跨时态精密切换”与复杂从句语序}
\begin{enumerate}
  \item \textbf{三大时态在历史政论中的精准分工}：\\
    \ding{172} \textbf{一般过去时（Datid）}：定位历史上已发生且结束的节点事件（如 \textit{i 1864 mistede Danmark..., i 1992 sagde vælgerne nej..., i 2022 stemte et flertal...}）。\\
    \ding{173} \textbf{现在完成时（Førnutid）}：连接过去决策与当代持续影响（如 \textit{Danmark har aldrig haft råd til..., I de senere år har krisen tvunget landet til...}）。\\
    \ding{174} \textbf{一般现在时（Nutid）}：阐述当今现实格局、结构性规律与未来追问（如 \textit{Spørgsmålet er..., Erfaringen viser, at...}）。
  \item \textbf{从句内部的否定词位置永远前移}：即使在 \textit{eftersom, idet, hvorvidt, hvormed} 引导的极正式从句中，否定词 \textit{ikke} 或副词 \textit{ofte / sjældent} 依然必须严格站在变位动词前面：\\
    \ding{51}\ \textit{...idet små stater \dabold{sjældent kan} forme deres egen sikkerhed alene.}（\ding{55}\ \textit{*...idet små stater kan sjældent forme...}）
\end{enumerate}
\end{grammatikbox}

\begin{grammatikbox}{高阶连词四大常见失误排雷（Top 4 Pitfalls）}
\begin{enumerate}
  \item \textbf{\ding{172} 误把 \textit{hvormed / idet} 当成 \textit{derfor} 并实施倒装}：\ding{55}\ \textit{*Landet mistede territorium, hvormed opstod en ny mentalitet.} $\rightarrow$ \ding{51}\ \textit{Landet mistede territorium, \dabold{hvormed} en ny mentalitet opstod.}（记住：\textit{hvormed, hvorved, idet, eftersom} 全是从属连词，从句内主语先于动词！）
  \item \textbf{\ding{173} 句首 \textit{Desuagtet} 忘记 V2 倒装或乱加逗号}：\ding{55}\ \textit{*Desuagtet, regeringen fulgte direktiverne.} $\rightarrow$ \ding{51}\ \textit{\dabold{Desuagtet} fulgte regeringen direktiverne.}
  \item \textbf{\ding{174} 混淆 \textit{i hvilken udstrækning}（在多大程度上）中的格与语序}：由于 \textit{udstrækning} 为共性单数名词（\textit{en udstrækning}），疑问限定词必须用 \dabold{hvilken}（\ding{55}\ \textit{*i hvilket udstrækning}），且后面从句不能倒装：\ding{51}\ \textit{i \dabold{hvilken} udstrækning vi vil sidde med ved bordet}。
  \item \textbf{\ding{175} 在 \textit{herunder} 后面画蛇添足加动词或 \textit{at}}：\textit{herunder} 用于直接在逗号后列举同位名词成分（\textit{flere forbehold, herunder euroforbeholdet}），不能接完整句子除非放在句首并倒装（\textit{Herunder spiller...}）。
\end{enumerate}
\end{grammatikbox}

% ---------------------------------------------------------------------
\ncdsubsec{五、Mønsterøvelser（政论高阶句法操练）}

\textbf{练习 A：C1 书面连词与历史政治核心词精准填空（每空填入最准确的 1 个丹麦语词）}
\begin{enumerate}
  \item Danmark blev i 1864 en sårbar småstat, \underline{\hspace{2.2cm}} en dyb forsigtighed slog rod i kulturen.
  \item Sammenholdet blev livsvigtigt, \underline{\hspace{2.2cm}} nationens overlevelse afhang af indre styrke.
  \item Landet valgte at deltage i det europæiske marked, \underline{\hspace{2.2cm}} åben handel sikrede velstanden.
  \item Eksporten hviler på flere stærke brancher, \underline{\hspace{2.2cm}} søfart og grøn energiteknologi.
  \item Vælgerne var uenige om, \underline{\hspace{2.2cm}} et tættere EU-samarbejde ville svække demokratiet.
  \item Debatten handler nu om, i hvilken \underline{\hspace{2.2cm}} et lille land bør dele suverænitet med andre.
  \item Befolkningen stemte nej i 1992; \underline{\hspace{2.2cm}} har Danmark ofte indført EU-regler meget hurtigt.
  \item Et flertal af vælgerne valgte i 2022 at \underline{\hspace{2.2cm}} det danske forsvarsforbehold.
\end{enumerate}

\newpage
\textbf{练习 B：高阶书面连接句纠错与升级（修正从句语序、V2 或用词错误，或按提示整合升级）}
\begin{enumerate}
  \item 【从句语序纠错】\textit{Småstater må ofte indgå alliancer, idet kan de sjældent klare sikkerheden alene.} $\rightarrow$
  \item 【句首倒装纠错】\textit{Desuagtet, Danmark har længe været en loyal partner i det europæiske miljøsamarbejde.} $\rightarrow$
  \item 【词形与语序纠错】\textit{Vi må undersøge, i hvilket udstrækning påvirker europæiske regler dansk erhvervsliv.} $\rightarrow$
  \item 【因果升级】用 \textit{eftersom} 合并两句：\textit{Danmark har få naturlige råstoffer. Landet er dybt afhængigt af global handel.} $\rightarrow$
  \item 【手段结果升级】用 \textit{, hvormed} 合并两句：\textit{Vælgerne fjernede forsvarsforbeholdet i 2022. Dermed fik Danmark en ny rolle i Europa.} $\rightarrow$
  \item 【归属列举升级】用 \textit{, herunder} 浓缩句子：\textit{Politikerne diskuterede flere danske forbehold, og de talte blandt andet om euroen og retsområdet.} $\rightarrow$
\end{enumerate}

\textbf{练习 C：中译丹（灵活调用本课高阶书面连词与历史政治词汇）}
\begin{enumerate}
  \item 丹麦依赖开放贸易，因为这个小国几乎没有自然资源。
  \item 选民们争论过，这是否会削弱国家的民主自决权。
  \item 关键问题在于，我们在多大程度上愿意让渡主权。
  \item 历史带来了疑虑，尽管如此，丹麦在实践中仍然忠实遵守欧洲规则。
  \item 该国出口多种商品，其中包括绿色能源和医药产品。
  \item 政府通过了新协议，从而加强了自身的国际影响力。
\end{enumerate}

% ---------------------------------------------------------------------
\ncdsubsec{六、PD3 / Studieprøven 实战微任务与参考答案}

\begin{emailbox}{Studieprøven / PD3 历史政论写作任务（220--260 ord）}
\textbf{Opgave}: Danmark beskrives ofte som et land, der både værner om sin nationale selvstændighed og samtidig er dybt afhængigt af omverdenen. Skriv en \textbf{sammenhængende historisk-politisk kommentar på 220--260 ord} om spændingsforholdet mellem småstatsmentalitet og globalt udsyn. Brug mindst \textbf{fem} af de syv avancerede forbindere fra lektionen (\textit{eftersom, idet, hvorvidt, desuagtet, herunder, hvormed, i hvilken udstrækning}) med korrekt ordstilling.
\end{emailbox}

\begin{answerbox}
\textbf{Før du læser 参考回答}
\begin{enumerate}
  \item \textit{Nederlaget i 1864 gjorde Danmark til en sårbar småstat, hvormed en dyb forsigtighed over for stormagter slog rod. Alligevel kan småstatsmentalitet og åbenhed eksistere side om side, fordi et lille land uden råstoffer har brug for fri handel og multilaterale alliancer for at sikre både velstand og sikkerhed.}
  \item \textit{Thomas Lind peger på det paradoks, at formelt at stå uden for EU-samarbejdet ikke sikrer reel selvbestemmelse, eftersom store beslutninger om handel, sikkerhed og grønne standarder alligevel træffes af nabolandene, så Danmark i praksis må følge direktiverne trods forbeholdene.}
\end{enumerate}

\textbf{Tekst B 阅读核对答案}
\begin{enumerate}
  \item \da{Forkert}（他指出形式上置身事外并不能保证真正的自决权，因为邻国仍会做出关键决定）
  \item \da{For ikke at isolere dansk erhvervsliv.}（为了不孤立丹麦的工商业界）
  \item \da{I hvilken udstrækning Danmark vil sidde med ved bordet, når reglerne skrives.}
\end{enumerate}

\textbf{练习 A 答案}
\begin{enumerate}
  \item \da{hvormed}（关系副词连词：由此，借此，回指前句变成小国的事实）
  \item \da{eftersom}（正式因果连词：因为，鉴于）
  \item \da{idet}（书面解释说明连词：因为，在于）
  \item \da{herunder}（副词：其中包括，后直接接名词短语）
  \item \da{hvorvidt}（正式间接疑问连词：究竟是否）
  \item \da{udstrækning}（固定搭配：\textit{i hvilken udstrækning} 在多大程度上）
  \item \da{desuagtet}（句间让步副词：尽管如此，分号后放句首触发 V2 倒装 \textit{har Danmark}）
  \item \da{ophæve}（动词原形：废除，取消，搭配保留条款）
\end{enumerate}

\textbf{练习 B 答案（参考）}
\begin{enumerate}
  \item \textit{Småstater må ofte indgå alliancer, idet de sjældent kan klare sikkerheden alene.}（\textit{idet} 引导从句，主语 \textit{de} 在前，副词 \textit{sjældent} 站在动词 \textit{kan} 前）
  \item \textit{Desuagtet har Danmark længe været en loyal partner i det europæiske miljøsamarbejde.}（去掉逗号，\textit{Desuagtet} 置句首立即触发 V2 倒装 \textit{har Danmark}）
  \item \textit{Vi må undersøge, i hvilken udstrækning europæiske regler påvirker dansk erhvervsliv.}（共性名词用 \textit{hvilken}，且间接疑问从句内部维持主语 + 动词不倒装）
  \item \textit{Danmark er dybt afhængigt af global handel, eftersom landet har få naturlige råstoffer.}（用正式因果连词 \textit{eftersom} 整合因果逻辑）
  \item \textit{Vælgerne fjernede forsvarsforbeholdet i 2022, hvormed Danmark fik en ny rolle i Europa.}（用 \textit{, hvormed} 顺滑承接结果，从句不倒装）
  \item \textit{Politikerne diskuterede flere danske forbehold, herunder euroen og retsområdet.}（用 \textit{herunder} 精炼替换冗长并列句）
\end{enumerate}

\textbf{练习 C 答案}
\begin{enumerate}
  \item \textit{Danmark er afhængigt af åben handel, eftersom det lille land næsten ingen råstoffer har.}
  \item \textit{Vælgerne diskuterede, hvorvidt dette ville svække landets demokratiske selvbestemmelse.}
  \item \textit{Det centrale spørgsmål er, i hvilken udstrækning vi vil afgive suverænitet.}
  \item \textit{Historien skabte skepsis; desuagtet følger Danmark i praksis loyalt de europæiske regler.}
  \item \textit{Landet eksporterer mange varer, herunder grøn energi og medicinske produkter.}
  \item \textit{Regeringen vedtog en ny aftale, hvormed den styrkede sin internationale indflydelse.}
\end{enumerate}

\textbf{PD3 / Studieprøven 高分范文（约 235 词）}\\
\textit{\textbf{Mellem forsigtighed og globalt ansvar: Småstatens europæiske balance}}

\textit{Gennem mere end halvandet århundrede har Danmarks forhold til Europa været præget af et dobbelt hensyn. Efter nederlaget i 1864 opstod en dyb \da{småstatsmentalitet}, \dabold{hvormed} forsigtighed over for stormagter og værn om national sammenhængskraft blev politiske kerneværdier. Landet lærte at opbygge stærke indre fællesskaber gennem folkehøjskoler og andelsbevægelse, \dabold{eftersom} en sårbar nation uden stor befolkning ikke kunne forlade sig på militær styrke.}

\textit{Alligevel har denne naturlige forsigtighed aldrig ført til varig økonomisk isolation. Tværtimod har dansk velstand siden efterkrigstiden hvilet på international åbenhed og frihandel, \dabold{idet} et lille hjemmemarked uden store råstoffer ikke kan bære en moderne velfærdsøkonomi alene. I dag bygger landets globale profil på specialiserede sektorer, \dabold{herunder} søfart, medicinalindustri og grøn energiteknologi, som alle forudsætter forudsigelige og fælles internationale spilleregler.}

\textit{Netop dette krydspres mellem national tryghed og global afhængighed forklarer den mangeårige debat om europæisk integration. Siden nej-stemmen i 1992 har danske vælgere gentagne gange spurgt, \dabold{hvorvidt} tættere samarbejde truer den lokale selvbestemmelse, og \dabold{i hvilken udstrækning} man bør \da{afgive} formel \da{suverænitet} til Bruxelles. \dabold{Desuagtet} viser nyere historie, at Danmark i praksis sjældent står udenfor, når afgørende europæiske standarder fastlægges. Ophævelsen af forsvarsforbeholdet i 2022 understregede yderligere, at reel sikkerhed og demokratisk indflydelse for et lille land ikke skabes i ensomhed.}

\textit{Fremtidens vigtigste opgave for dansk \da{udenrigspolitik} er derfor ikke at vælge mellem nære rødder og \da{globalt udsyn}, men at forene begge dele gennem aktivt, pragmatisk og konstruktivt europæisk engagement.}

\small\textbf{名师点评：}这篇约 235 词的范文堪称历史政论文的语法教科书示范：全部七大目标书面连接工具（\textit{hvormed, eftersom, idet, herunder, hvorvidt, i hvilken udstrækning, desuagtet}）尽数自然嵌入论证肌理，从句内语序与句首倒装无一疏漏；同时精准示范了 1864 史实过去时、持续影响的现在完成时与当下结论的现在时三层切换。
\end{answerbox}
